\documentclass[10pt,a4paper]{amsart}
\usepackage{amsmath,amssymb,amsthm,graphicx}
% Retained minimal preamble for the Stacks Project chapter more-morphisms.tex (More on Morphisms).
% Verbatim excerpt of the chapter preamble: Xy-pic setup (native preamble lines 14-18) and the native
% theorem declarations with the mathematical macros (native preamble lines 154-217). Copied mechanically;
% nothing added, removed or redefined. The wrapper class and the amsmath/amssymb/amsthm/graphicx base
% packages are supplied by the assembler, so the class/externaldocument part of the native preamble is omitted.
% For commutative diagrams we use Xy-pic
\usepackage[all]{xy}

% We use 2cell for 2-commutative diagrams.
\xyoption{2cell}
% Theorem environments.
%
\theoremstyle{plain}
\newtheorem{theorem}[subsection]{Theorem}
\newtheorem{proposition}[subsection]{Proposition}
\newtheorem{lemma}[subsection]{Lemma}

\theoremstyle{definition}
\newtheorem{definition}[subsection]{Definition}
\newtheorem{example}[subsection]{Example}
\newtheorem{exercise}[subsection]{Exercise}
\newtheorem{situation}[subsection]{Situation}

\theoremstyle{remark}
\newtheorem{remark}[subsection]{Remark}
\newtheorem{remarks}[subsection]{Remarks}

\numberwithin{equation}{subsection}

% Macros
%
\def\lim{\mathop{\mathrm{lim}}\nolimits}
\def\colim{\mathop{\mathrm{colim}}\nolimits}
\def\Spec{\mathop{\mathrm{Spec}}}
\def\Hom{\mathop{\mathrm{Hom}}\nolimits}
\def\Ext{\mathop{\mathrm{Ext}}\nolimits}
\def\SheafHom{\mathop{\mathcal{H}\!\mathit{om}}\nolimits}
\def\SheafExt{\mathop{\mathcal{E}\!\mathit{xt}}\nolimits}
\def\Sch{\mathit{Sch}}
\def\Mor{\mathop{\mathrm{Mor}}\nolimits}
\def\Ob{\mathop{\mathrm{Ob}}\nolimits}
\def\Sh{\mathop{\mathit{Sh}}\nolimits}
\def\NL{\mathop{N\!L}\nolimits}
\def\CH{\mathop{\mathrm{CH}}\nolimits}
\def\proetale{{pro\text{-}\acute{e}tale}}
\def\etale{{\acute{e}tale}}
\def\QCoh{\mathit{QCoh}}
\def\Ker{\mathop{\mathrm{Ker}}}
\def\Im{\mathop{\mathrm{Im}}}
\def\Coker{\mathop{\mathrm{Coker}}}
\def\Coim{\mathop{\mathrm{Coim}}}

% Boxtimes
%
\DeclareMathSymbol{\boxtimes}{\mathbin}{AMSa}{"02}

%
% Macros for moduli stacks/spaces
%
\def\QCohstack{\mathcal{QC}\!\mathit{oh}}
\def\Cohstack{\mathcal{C}\!\mathit{oh}}
\def\Spacesstack{\mathcal{S}\!\mathit{paces}}
\def\Quotfunctor{\mathrm{Quot}}
\def\Hilbfunctor{\mathrm{Hilb}}
\def\Curvesstack{\mathcal{C}\!\mathit{urves}}
\def\Polarizedstack{\mathcal{P}\!\mathit{olarized}}
\def\Complexesstack{\mathcal{C}\!\mathit{omplexes}}
% \Pic is the operator that assigns to X its picard group, usage \Pic(X)
% \Picardstack_{X/B} denotes the Picard stack of X over B
% \Picardfunctor_{X/B} denotes the Picard functor of X over B
\def\Pic{\mathop{\mathrm{Pic}}\nolimits}
\def\Picardstack{\mathcal{P}\!\mathit{ic}}
\def\Picardfunctor{\mathrm{Pic}}
\def\Deformationcategory{\mathcal{D}\!\mathit{ef}}
\usepackage{cleveref}
\crefname{originalreference}{original reference}{original references}
\Crefname{originalreference}{Original reference}{Original references}
\makeatletter
\newlabel{limits-lemma-descend-finite-presentation}{{01ZM}{source}{}{sourcepointer-1}{}}
\newlabel{limits-lemma-descend-finite-presentation@cref}{{[originalreference][0][]01ZM}{[1][0][]source}{}{}{}}
\makeatother
\crefname{originalreference}{original reference}{original references}
\Crefname{originalreference}{Original reference}{Original references}
\makeatletter
\newlabel{limits-lemma-descend-modules-finite-presentation}{{01ZR}{source}{}{sourcepointer-2}{}}
\newlabel{limits-lemma-descend-modules-finite-presentation@cref}{{[originalreference][0][]01ZR}{[1][0][]source}{}{}{}}
\makeatother
\crefname{originalreference}{original reference}{original references}
\Crefname{originalreference}{Original reference}{Original references}
\makeatletter
\newlabel{limits-proposition-approximate}{{01ZA}{source}{}{sourcepointer-3}{}}
\newlabel{limits-proposition-approximate@cref}{{[originalreference][0][]01ZA}{[1][0][]source}{}{}{}}
\makeatother
\crefname{originalreference}{original reference}{original references}
\Crefname{originalreference}{Original reference}{Original references}
\makeatletter
\newlabel{morphisms-proposition-generic-flatness-reduced}{{052B}{source}{}{sourcepointer-4}{}}
\newlabel{morphisms-proposition-generic-flatness-reduced@cref}{{[originalreference][0][]052B}{[1][0][]source}{}{}{}}
\makeatother
\title{Generic flatness stratification for a finitely presented module}
\author{The Stacks Project Authors}
\date{stacks-project:more-morphisms.tex:local-snapshot-20261011; GNU FDL 1.2 or later}
\begin{document}
\maketitle
\noindent\textit{Editorial note.} Counted span is the author's own proof environment, source lines 15579-15612, immediately following the statement at 15565-15577.\par\medskip
\noindent\textit{Editorial note.} Four cross-chapter pointers (Limits, Morphisms) are recorded in assembly.external\_references.\par\medskip
\noindent\textit{Editorial note.} The statement's descending chain with its three conditions is retained verbatim in the statement span.\par\medskip
\noindent\textit{Allowed standard prerequisites (editorial).}\begin{itemize}
\item Generic flatness over a reduced base. Morphisms, proposition morphisms-proposition-generic-flatness-reduced, cited at proof lines 15602-15604. Supplies the dense open U\_0 over which the pullback of F is flat, which drives the descending chain.
\item Approximation of a finitely presented morphism. Limits, proposition limits-proposition-approximate with the descent lemmas cited in the proof. Reduces the general statement to a Noetherian situation where the chain terminates.
\item Descending finite presentation and finite modules. Limits, lemmas limits-lemma-descend-finite-presentation and limits-lemma-descend-modules-finite-presentation, cited in the proof. Keeps the finiteness hypotheses along the approximation.
\item Descending chains of closed subsets in a Noetherian space. used at proof lines 15609-15611. Gives termination of the constructed chain, producing the required t.
\end{itemize}\medskip
\ifdefined\hypertarget\hypertarget{sourcepointer-1}{}\fi
\noindent\textit{Source reference.} Original reference 01ZM: focused-seeds/geometry/S9-morph-generic-flatness-stratification/ref-tags.txt, source lines 5--5. Navigation pointer only: the counted proof cites this statement by \textbackslash{}ref as a standing input; the deduction the counted proof has to make is carried out inside the counted span, and no part of that deduction is delegated to this pointer.\par\medskip
\ifdefined\hypertarget\hypertarget{sourcepointer-2}{}\fi
\noindent\textit{Source reference.} Original reference 01ZR: focused-seeds/geometry/S9-morph-generic-flatness-stratification/ref-tags.txt, source lines 6--6. Navigation pointer only: the counted proof cites this statement by \textbackslash{}ref as a standing input; the deduction the counted proof has to make is carried out inside the counted span, and no part of that deduction is delegated to this pointer.\par\medskip
\ifdefined\hypertarget\hypertarget{sourcepointer-3}{}\fi
\noindent\textit{Source reference.} Original reference 01ZA: focused-seeds/geometry/S9-morph-generic-flatness-stratification/ref-tags.txt, source lines 7--7. Navigation pointer only: the counted proof cites this statement by \textbackslash{}ref as a standing input; the deduction the counted proof has to make is carried out inside the counted span, and no part of that deduction is delegated to this pointer.\par\medskip
\ifdefined\hypertarget\hypertarget{sourcepointer-4}{}\fi
\noindent\textit{Source reference.} Original reference 052B: focused-seeds/geometry/S9-morph-generic-flatness-stratification/ref-tags.txt, source lines 8--8. Navigation pointer only: the counted proof cites this statement by \textbackslash{}ref as a standing input; the deduction the counted proof has to make is carried out inside the counted span, and no part of that deduction is delegated to this pointer.\par\medskip
\begin{lemma}[Generic flatness stratification]
\label{lemma-generic-flatness-stratification}
Let $f : X \to S$ be a morphism of finite presentation between quasi-compact
and quasi-separated schemes. Let $\mathcal{F}$ be an $\mathcal{O}_X$-module
of finite presentation. Then there exists a $t \geq 0$ and closed
subschemes
$$
S \supset S_0 \supset S_1 \supset \ldots \supset S_t = \emptyset
$$
such that $S_i \to S$ is defined by a finite type ideal sheaf,
$S_0 \subset S$ is a thickening, and $\mathcal{F}$ pulled back to
$X \times_S (S_i \setminus S_{i + 1})$ is flat over $S_i \setminus S_{i + 1}$.
\end{lemma}

\begin{proof}
We can find a cartesian diagram
$$
\xymatrix{
X \ar[d] \ar[r] & X_0 \ar[d] \\
S \ar[r] & S_0
}
$$
and a finitely presented $\mathcal{O}_{X_0}$-module $\mathcal{F}_0$
which pulls back to $\mathcal{F}$ such that $X_0$ and $S_0$ are of
finite type over $\mathbf{Z}$. See
Limits, Proposition \ref{limits-proposition-approximate} and
Lemmas \ref{limits-lemma-descend-finite-presentation} and
\ref{limits-lemma-descend-modules-finite-presentation}.
Thus we may assume $X$ and $S$ are of finite type over $\mathbf{Z}$
and $\mathcal{F}$ is a coherent $\mathcal{O}_X$-module.

\medskip\noindent
Assume $X$ and $S$ are of finite type over $\mathbf{Z}$
and $\mathcal{F}$ is a coherent $\mathcal{O}_X$-module.
In this case every quasi-coherent ideal is of finite type, hence
we do not have to check the condition that $S_i$ is cut out
by a finite type ideal. Set $S_0 = S_{red}$ equal to the reduction of $S$.
By generic flatness as stated in Morphisms, Proposition
\ref{morphisms-proposition-generic-flatness-reduced}
there is a dense open $U_0 \subset S_0$ such that $\mathcal{F}$
pulled back to $X \times_S U_0$ is flat over $U_0$.
Let $S_1 \subset S_0$ be the reduced closed subscheme whose
underlying closed subset is $S \setminus U_0$. We continue in this
way, provided $S_1 \not = \emptyset$, to find
$S_0 \supset S_1 \supset \ldots$. Because $S$
is Noetherian any descending chain of closed subsets stabilizes
hence we see that $S_t = \emptyset$ for some $t \geq 0$.
\end{proof}

\end{document}
